% TOP_PROBLEM: {"id":30007556,"problem_number":"LOCAL-30007556","title":"Economic models of costly computation","category_id":21,"category":{"id":21,"name":"economics","display_name":"Economics","description":"Open economic theory statements, published research agendas, and proposed scoped specifications; question provenance and historical priority are qualified per record.","slug":"economics","order_index":21,"created_at":"2026-10-08T00:00:00Z"},"status":"research_question","external_url":null,"legacy_ids":[],"published":false,"statement":"Identify computational-cost and algorithm-repertoire models that jointly predict choices and decision times across different economic problem families.","economics":{"original_id":45,"field":"Economic theory, equilibrium and social choice","kind":"design","formalization_basis":"adapted_research_question","source_status":"research_agenda","priority_status":"earliest_found","priority_note":"This is the earliest source in which the relevant agenda was verified during this audit; absolute priority has not been established.","earliest_verified_reference":{"authors":["Geoffroy de Clippel","Kareen Rozen"],"title":"Bounded Rationality in Choice Theory: A Survey","year":2023,"url":"https://geoffroydeclippel.net/Working%20Papers%20PDFs/JEL.pdf","doi":"10.1257/jel.20231592","locator":"§6.4 Searching for common ground; §7 Looking Forward, PDF pp.51–54","evidence_summary":"The revised working paper identifies welfare comparisons, additional process data, and computational or procedural constraints as continuing research directions."},"current_reference":{"authors":["Geoffroy de Clippel","Kareen Rozen"],"title":"Bounded Rationality in Choice Theory: A Survey","year":2023,"url":"https://geoffroydeclippel.net/Working%20Papers%20PDFs/JEL.pdf","doi":"10.1257/jel.20231592","locator":"§6.4 Searching for common ground; §7 Looking Forward, PDF pp.51–54","evidence_summary":"The revised working paper identifies welfare comparisons, additional process data, and computational or procedural constraints as continuing research directions."},"formalization_note":"The computational repertoire and validation criterion are introduced to make the survey’s broad procedural-model agenda reviewable."},"statusQualification":"Published question or research agenda verified at the cited locator. The mathematical specification is adapted for this catalogue and includes additional assumptions; source publication does not establish first-ever priority or certify this exact model remains unsolved. The computational repertoire and validation criterion are introduced to make the survey’s broad procedural-model agenda reviewable.","statusReviewedAt":"2026-10-08"}
% FIRST_POSTED: 2026-10-08
% ENTRY_KIND: statement_only
% !TeX program = lualatex
\documentclass[11pt]{article}
\usepackage{fontspec}
\usepackage[a4paper,margin=25mm]{geometry}
\usepackage{amsmath,amssymb,mathtools}
\usepackage{xurl}
\usepackage[hidelinks]{hyperref}
\newcommand{\sref}[2]{\hyperlink{source:#1:#2}{[S#2]}}
\setlength{\emergencystretch}{3em}
\begin{document}
% problemId:30007556
\section*{Economic models of costly computation}\label{30007556}
\subsection{Definitions and mathematical statement}
Consider economic decision problems \(d\) drawn from a distribution \(\mathcal D\). Each problem specifies a finite feasible set \(X_d\), payoff function \(u_d:X_d\to[0,1]\), and an information representation accessible through queries. A decision algorithm \(A\) maps query histories into further queries or a final choice \(x\). Let \(C(A,d)\) be its expected number of elementary operations, and let \(\kappa\geq0\) be an individual's cost per operation. A computational-choice model selects an algorithm from a stated repertoire \(\mathcal A\) to maximize
\[
 V(A,d,\kappa)=E[u_d(A(d))]-\kappa C(A,d).
\]
The target is to construct and identify a parsimonious repertoire and cost specification that jointly predict choices, response times, and the effects of experimentally varied information complexity. Estimate \(\kappa\) and repertoire heterogeneity using randomized changes in representation while holding material payoffs fixed. Determine restrictions under which algorithmic cost can be distinguished from preferences over outcomes and statistical choice noise. Require successful prediction on new problem families and an explicit account of how primitive computational operations map into measured time. A solution may establish failure of portable parameterization and identify the additional primitives required. This is an empirical modeling and identification agenda; optimizing a fully known finite repertoire by itself would not resolve the portability problem.

\par\textbf{Formulation and scope.} The computational repertoire and validation criterion are introduced to make the survey’s broad procedural-model agenda reviewable.

\par\textbf{Open-question provenance.} An explicit question or research agenda was verified at the source locator. This does not by itself certify that every later version remains unresolved \sref{30007556}{1}.

\par\textbf{Historical priority.} This is the earliest source in which the relevant agenda was verified during this audit; absolute priority has not been established.

\subsection{Short English statement}
Identify computational-cost and algorithm-repertoire models that jointly predict choices and decision times across different economic problem families.

\subsection{Sources}
\begin{itemize}\raggedright
\item[\textnormal{[S1]}]\hypertarget{source:30007556:1}{} Geoffroy de Clippel, Kareen Rozen. 2023. Bounded Rationality in Choice Theory: A Survey. §6.4 Searching for common ground; §7 Looking Forward, PDF pp.51–54.
\url{https://geoffroydeclippel.net/Working%20Papers%20PDFs/JEL.pdf}.
DOI: 10.1257/jel.20231592.
{\small\textit{Earliest verified question/agenda reference; historical priority qualified below. The revised working paper identifies welfare comparisons, additional process data, and computational or procedural constraints as continuing research directions.}}
\end{itemize}
\end{document}
